\chapter{Morita equivalence for full corners}
\label{ch:rm03-morita}

We use the right-module convention in this chapter. The progenerator
form of Morita equivalence and its converse are proved in
\cite[Ch.~3, ALG-03-C02-S05]{p103}. In that theorem a finitely generated
projective generator \(P_R\) gives inverse equivalences by tensoring
with \(\operatorname{Hom}_R(P,R)\) and with \(P\) on the indicated
sides. Here we identify the generator condition and the two functors
for a corner
\cite[Lec.~7]{mitmorita}.

\section{The full-idempotent criterion}

\begin{corollary}[RM03-C05.01. Morita equivalence from a full idempotent]
\label{cor:rm03-full-corner}
Let \(R\) be a unital ring and \(e\in R\) satisfy \(e^2=e\). The
right module \(eR\) is finitely generated projective, and the map
\[
   eRe\longrightarrow\operatorname{End}_R(eR),\qquad
   a\longmapsto (x\longmapsto ax)
\]
is an isomorphism of unital rings when \(eRe\) has identity \(e\).
The module \(eR\) is a generator of right \(R\)-modules if and only if
\(e\) is full, meaning that \(ReR=R\). In this case there are inverse
equivalences of right-module categories
\[
   \mathrm{Mod}\text{-}R \longrightarrow \mathrm{Mod}\text{-}eRe,
   \quad M\longmapsto Me,
   \qquad
   N\longmapsto N\otimes_{eRe}eR
\]
with the second functor directed back to \(\mathrm{Mod}\text{-}R\).
\end{corollary}
\begin{proof}
The decomposition \(R_R=eR\oplus(1-e)R\) follows by multiplying on
the left by \(e\) and \(1-e\). Therefore \(eR\) is projective and is
generated by \(e\). Every \(R\)-endomorphism \(f\) of \(eR\) has
\(f(e)=f(e)e\in eRe\) and satisfies \(f(er)=f(e)r\). Evaluation at
\(e\) is inverse to the displayed left multiplication map. Products
correspond to composition and the identity endomorphism to \(e\).

If \(e\) is full, a finite expression
\(1=\sum_{i=1}^m r_i e s_i\) gives an \(R\)-linear surjection
\[
    (eR)^m\longrightarrow R_R,\qquad
    (x_i)_{i=1}^m\longmapsto\sum_i r_i x_i,
\]
since \((e s_i)_i\) maps to \(1\). Every right module is a quotient
of a direct sum of copies of \(R_R\), so every right module is a
quotient of a direct sum of copies of \(eR\). Conversely, if \(eR\)
is a generator, take a surjection from a direct sum of copies of
\(eR\) onto \(R_R\) and a finite-support preimage of \(1\). Each
component \(g_i:eR\to R\) has \(g_i(e)=g_i(e)e\) and takes \(e s\)
to \(g_i(e)e s\in ReR\). Thus \(1\in ReR\), which proves fullness.

Evaluation at \(e\) identifies \(\operatorname{Hom}_R(eR,R)\)
with \(Re\). The inverse of evaluation sends \(b\in Re\) to
\(x\mapsto bx\) on \(eR\). Multiplication on the left by \(R\)
and precomposition by \(a\in eRe\) become the usual left \(R\)- and
right \(eRe\)-actions on \(Re\), respectively. Thus the first
functor of the progenerator theorem is \(M\otimes_R Re\).
The natural right \(eRe\)-linear map \(m\otimes b\mapsto mb\)
identifies this tensor product with \(Me\). Its inverse sends
\(u\in Me\) to \(u\otimes e\). One composite is \(ue=u\), and
the other is \(mb\otimes e=m\otimes be=m\otimes b\).
The inverse functor in the theorem is \(N\mapsto N\otimes_{eRe}eR\).
\end{proof}

Under these identifications the counit is the multiplication map
\(Me\otimes_{eRe}eR\to M\). One can read off its inverse from a full
expression \(1=\sum_i r_i e s_i\). It sends \(m\) to
\(\sum_i (m r_i e)\otimes(e s_i)\). Multiplying gives \(m\) back.
For \(m'\in Me\) and \(x\in eR\), the element \(x r_i e\) belongs to
\(eRe\). Balancing moves this element across the tensor sign and
gives the other composition
\(\sum_i m'\otimes x r_i e s_i=m'\otimes x\).

\begin{example}[RM03-C05.02. A matrix corner and a non-full corner]
\label{ex:rm03-corners}
For \(n\geq 1\), \(e_{11}\) is full in \(M_n(R)\), with corner
isomorphic to \(R\). In the product of two nonzero rings \(A\times B\),
the corner of \(e=(1,0)\) is isomorphic to \(A\), and the corner
functor \(M\mapsto Me\) kills the nonzero right module \(B\).
\end{example}
\begin{proof}
The matrix units satisfy \(\sum_i e_{i1}e_{11}e_{1i}=1_{M_n(R)}\).
The corner is the position \((1,1)\), and its product is the product
of \(R\). Thus Corollary~\ref{cor:rm03-full-corner} gives the matrix
equivalence. For the product ring the ideal generated by \((1,0)\)
is \(A\times 0\), and right multiplication by \((1,0)\) on the
module \(B\) obtained from projection to \(B\) is zero.
\end{proof}
